% Lösungen Einheit 1
\einheitenkopf*{Einheit 1 von 3 · Kreisumfang}
\erg{1}{a) Radius \quad b) Durchmesser \quad c) Radius \quad d) Durchmesser \quad e) Radius}
\erg{2}{a) Umfang \quad b) Flächeninhalt \quad c) Umfang \quad d) Flächeninhalt \quad e) Umfang}
\erg{3}{a) Zirkel 2,5\,cm weit öffnen, Einstich in M \quad b) $r = 7 : 2 = 3{,}5$\,cm, Zirkel 3,5\,cm weit öffnen \quad c) Radius: Strecke von M zum Rand; Durchmesser: Strecke durch M von Rand zu Rand (7\,cm lang)}
\erg{4}{a) $d = 12$\,cm \quad b) $r = 9$\,cm \quad c) $r = 5{,}5$\,cm \quad d) $d = 7$\,m \quad e) $r = 0{,}65$\,m \quad f) $r = 50$\,cm}
\erg{5}{a) 3,16 \quad b) 3,13 \quad c) 3,15 \quad d) 3,14 \quad e) ungefähr 3,14-mal (gut dreimal)}
\erg{6}{a) 9,4\,cm \quad b) 22,0\,cm \quad c) 25,1\,m \quad d) 37,7\,cm \quad e) 56,5\,cm \quad f) 8,2\,cm \quad g) 4,7\,m}
\erg{7}{a) 14,0\,cm \quad b) 31,8\,m \quad c) $d = 50 : \pi \approx 15{,}92$\,cm, $r \approx 8{,}0$\,cm \quad d) $b = \pi \cdot 9 : 2 \approx 14{,}1$\,cm \quad e) $14{,}14 + 9 \approx 23{,}1$\,cm \quad f) $u = 2 \cdot \pi \cdot 4{,}3 \approx 27{,}0$\,cm; das Blech muss etwa 27,0\,cm lang sein}
\erg{8}{a) Tom setzt den Durchmesser 14\,cm für den Radius ein. Richtig: $u = \pi \cdot 14\,\text{cm} \approx 44{,}0$\,cm (oder $r = 7$\,cm, $u = 2 \cdot \pi \cdot 7$) \quad b) 81,7\,m}
\erg{9}{a) Der Umfang verdoppelt sich, weil $u = \pi \cdot d$: doppelter Durchmesser, doppelter Umfang. \quad b) Nein: $u : d$ ist bei jedem Kreis dieselbe Zahl $\pi$, egal wie groß er ist. \quad c) Länger: Der Weg ist ein Umfang, also $\pi \cdot 1$\,m, und $\pi$ ist etwas größer als 3.}
\erg{10}{a) $u = \pi \cdot 70\,\text{cm} \approx 219{,}9$\,cm \quad b) $500 \cdot 219{,}9$\,cm $\approx 109\,956$\,cm $\approx 1099{,}6$\,m, also rund 1100\,m \quad c) $2500\,\text{m} : 2{,}199\,\text{m} \approx 1137$ Umdrehungen}
